% Generated draft: requires mathematical and source review.
\documentclass[11pt]{article}
\usepackage[margin=25mm]{geometry}
\usepackage{fontspec}
\setmainfont{Latin Modern Roman}
\usepackage{amsmath,amssymb,mathtools,unicode-math}
\setmathfont{Latin Modern Math}
\usepackage{enumitem,needspace,xurl,hyperref}
\setlength{\parindent}{0pt}
\setlength{\parskip}{5pt}
\setlength{\emergencystretch}{3em}
\newcommand{\N}{\mathbb{N}}
\newcommand{\Z}{\mathbb{Z}}
\newcommand{\Q}{\mathbb{Q}}
\newcommand{\R}{\mathbb{R}}
\newcommand{\C}{\mathbb{C}}
\newcommand{\F}{\mathbb{F}}
\newcommand{\A}{\mathbb{A}}
\newcommand{\PP}{\mathbb{P}}
\newcommand{\E}{\mathbb{E}}
\newcommand{\eps}{\varepsilon}
\newcommand{\dd}{\,\mathrm{d}}
\begin{document}
\section*{UnsolvedMath ID: 701104}
\textbf{Source problem code:} catalogue entry
\subsection{Definitions and mathematical statement}
Let \(G\) be a finitely generated group and \(S\) a finite symmetric generating set. Define its word-growth function by
\[
\gamma_{G,S}(n)=\bigl|\{g\in G:\lvert g\rvert_S\le n\}\bigr|,\qquad n\in\Z_{\ge0},
\]
where \(\lvert g\rvert_S\) is the minimum length of a word in \(S\) representing \(g\). For functions \(f,g:\Z_{\ge0}\to\R_{>0}\), write \(f\preccurlyeq g\) if there exist positive integers \(A,B\) and \(n_0\in\Z_{\ge0}\) such that \(f(n)\le A g(Bn)\) for every \(n\ge n_0\); write \(f\asymp g\) if both \(f\preccurlyeq g\) and \(g\preccurlyeq f\). The equivalence class of \(\gamma_{G,S}\) is independent of the finite generating set.

For the branch-group terminology, use the following standard convention. A rooted spherically homogeneous tree \(T\) has a root and, for each \(k\ge0\), every vertex at distance \(k\) from the root has the same finite number \(d_k\ge2\) of children. Let \(L_k\) be its \(k\)-th level. For an action of \(G\) by root-preserving tree automorphisms and a vertex \(v\), the rigid stabilizer \(\operatorname{rist}_G(v)\) consists of elements fixing every vertex outside the subtree rooted at \(v\). Put
\[
\operatorname{rist}_G(k)=\bigl\langle\operatorname{rist}_G(v):v\in L_k\bigr\rangle.
\]
A group is a branch group if it admits such a faithful action, transitive on every level, with \([G:\operatorname{rist}_G(k)]<\infty\) for every \(k\ge0\).

Under the interpretation that ``rate of growth is \(e^{\sqrt n}\)'' means growth equivalence \(\asymp\), the two questions are:
\begin{enumerate}
\item[(a)] Does there exist a finitely generated branch group \(G\) such that
\[
\gamma_{G,S}(n)\asymp \exp(\sqrt n)?
\]
\item[(b)] Does there exist a finitely generated group \(G\) such that
\[
\gamma_{G,S}(n)\asymp \exp(\sqrt n)?
\]
\end{enumerate}
In each question, \(S\) may be any finite symmetric generating set. The supplied record does not define its growth-comparison convention or its branch-group convention; the definitions above make explicit the interpretation used in this draft. In particular, the source does not specify literal equality, asymptotic ratio one, or two-sided constant multiples of \(\exp(\sqrt n)\). It also states that the question is related to Kourovka Notebook Problem 9.9, whose statement is not supplied.
\subsection{Short English statement}
Does a finitely generated branch group have word growth equivalent to \(\exp(\sqrt n)\)? Does any finitely generated group have that growth? This draft uses the stated coarse growth equivalence; the source leaves that convention implicit. The supplied registry labels the problem open, without independent verification here.
\subsection{Sources}
\begin{itemize}
\item Local UnsolvedMath registry: permanent ID \(701104\), source code \texttt{catalogue entry}, title ``Kourovka Notebook Problem 15.16.'' The supplied record labels it \texttt{open}; this is supplied source metadata, not an independently verified mathematical assessment.
\item \emph{The Kourovka Notebook (Unsolved Problems in Group Theory)}, Problem 15.16, originating in the 15th issue (2002); the supplied background also identifies the 21st-issue edition (2026), \texttt{arXiv:1401.0300v48}. The collection description names E. I. Khukhro and V. D. Mazurov as editors. This reference is supplied attribution; the original text has not been independently checked here.
\item Proposed in 2002 by L. Bartholdi, R. I. Grigorchuk, and a third proposer whose name is corrupted in the supplied record. The proposer field ends with the fragment \texttt{Z. \{\textbackslash v S}; the background separately contains \texttt{uni\textbackslash 'k\}}. The full name cannot be transcribed reliably from these fields.
\item The supplied statement refers to Kourovka Notebook Problem 9.9 as related, without supplying its statement or explaining the relationship. No usable discussion or further literature citation accompanies the corrupted discussion field.
\end{itemize}
\end{document}
